\documentclass[10pt,a4paper]{amsart}
\usepackage{amsmath,amssymb,amsthm,graphicx}
% For commutative diagrams we use Xy-pic
\usepackage[all]{xy}

% We use 2cell for 2-commutative diagrams.
\xyoption{2cell}
% Theorem environments.
%
\theoremstyle{plain}
\newtheorem{theorem}[subsection]{Theorem}
\newtheorem{proposition}[subsection]{Proposition}
\newtheorem{lemma}[subsection]{Lemma}

\theoremstyle{definition}
\newtheorem{definition}[subsection]{Definition}
\newtheorem{example}[subsection]{Example}
\newtheorem{exercise}[subsection]{Exercise}
\newtheorem{situation}[subsection]{Situation}

\theoremstyle{remark}
\newtheorem{remark}[subsection]{Remark}
\newtheorem{remarks}[subsection]{Remarks}

\numberwithin{equation}{subsection}

% Macros
%
\def\lim{\mathop{\mathrm{lim}}\nolimits}
\def\colim{\mathop{\mathrm{colim}}\nolimits}
\def\Spec{\mathop{\mathrm{Spec}}}
\def\Hom{\mathop{\mathrm{Hom}}\nolimits}
\def\Ext{\mathop{\mathrm{Ext}}\nolimits}
\def\SheafHom{\mathop{\mathcal{H}\!\mathit{om}}\nolimits}
\def\SheafExt{\mathop{\mathcal{E}\!\mathit{xt}}\nolimits}
\def\Sch{\mathit{Sch}}
\def\Mor{\mathop{\mathrm{Mor}}\nolimits}
\def\Ob{\mathop{\mathrm{Ob}}\nolimits}
\def\Sh{\mathop{\mathit{Sh}}\nolimits}
\def\NL{\mathop{N\!L}\nolimits}
\def\CH{\mathop{\mathrm{CH}}\nolimits}
\def\proetale{{pro\text{-}\acute{e}tale}}
\def\etale{{\acute{e}tale}}
\def\QCoh{\mathit{QCoh}}
\def\Ker{\mathop{\mathrm{Ker}}}
\def\Im{\mathop{\mathrm{Im}}}
\def\Coker{\mathop{\mathrm{Coker}}}
\def\Coim{\mathop{\mathrm{Coim}}}

% Boxtimes
%
\DeclareMathSymbol{\boxtimes}{\mathbin}{AMSa}{"02}

%
% Macros for moduli stacks/spaces
%
\def\QCohstack{\mathcal{QC}\!\mathit{oh}}
\def\Cohstack{\mathcal{C}\!\mathit{oh}}
\def\Spacesstack{\mathcal{S}\!\mathit{paces}}
\def\Quotfunctor{\mathrm{Quot}}
\def\Hilbfunctor{\mathrm{Hilb}}
\def\Curvesstack{\mathcal{C}\!\mathit{urves}}
\def\Polarizedstack{\mathcal{P}\!\mathit{olarized}}
\def\Complexesstack{\mathcal{C}\!\mathit{omplexes}}
% \Pic is the operator that assigns to X its picard group, usage \Pic(X)
% \Picardstack_{X/B} denotes the Picard stack of X over B
% \Picardfunctor_{X/B} denotes the Picard functor of X over B
\def\Pic{\mathop{\mathrm{Pic}}\nolimits}
\def\Picardstack{\mathcal{P}\!\mathit{ic}}
\def\Picardfunctor{\mathrm{Pic}}
\def\Deformationcategory{\mathcal{D}\!\mathit{ef}}
\usepackage{cleveref}
\crefname{originalreference}{original reference}{original references}
\Crefname{originalreference}{Original reference}{Original references}
\makeatletter
\newlabel{lemma-h0-topological}{{09ZG}{source}{}{sourcepointer-1}{}}
\newlabel{lemma-h0-topological@cref}{{[originalreference][0][]09ZG}{[1][0][]source}{}{}{}}
\makeatother
\crefname{originallemma}{original lemma}{original lemmas}
\Crefname{originallemma}{Original lemma}{Original lemmas}
\makeatletter
\newlabel{more-algebra-lemma-integral-over-henselian-pair}{{09XK}{source}{}{sourcepointer-2}{}}
\newlabel{more-algebra-lemma-integral-over-henselian-pair@cref}{{[originallemma][0][]09XK}{[1][0][]source}{}{}{}}
\makeatother
\crefname{originallemma}{original lemma}{original lemmas}
\Crefname{originallemma}{Original lemma}{Original lemmas}
\makeatletter
\newlabel{more-algebra-lemma-irreducible-henselian-pair-connected}{{09Y6}{source}{}{sourcepointer-3}{}}
\newlabel{more-algebra-lemma-irreducible-henselian-pair-connected@cref}{{[originallemma][0][]09Y6}{[1][0][]source}{}{}{}}
\makeatother
\crefname{originaltheorem}{original theorem}{original theorems}
\Crefname{originaltheorem}{Original theorem}{Original theorems}
\makeatletter
\newlabel{more-morphisms-theorem-stein-factorization-general}{{03H2}{source}{}{sourcepointer-4}{}}
\newlabel{more-morphisms-theorem-stein-factorization-general@cref}{{[originaltheorem][0][]03H2}{[1][0][]source}{}{}{}}
\makeatother
\crefname{originallemma}{original lemma}{original lemmas}
\Crefname{originallemma}{Original lemma}{Original lemmas}
\makeatletter
\newlabel{properties-lemma-quasi-compact-quasi-separated-spectral}{{094L}{source}{}{sourcepointer-5}{}}
\newlabel{properties-lemma-quasi-compact-quasi-separated-spectral@cref}{{[originallemma][0][]094L}{[1][0][]source}{}{}{}}
\makeatother
\crefname{originallemma}{original lemma}{original lemmas}
\Crefname{originallemma}{Original lemma}{Original lemmas}
\makeatletter
\newlabel{topology-lemma-connected-fibres-quotient-topology-connected-components}{{0377}{source}{}{sourcepointer-6}{}}
\newlabel{topology-lemma-connected-fibres-quotient-topology-connected-components@cref}{{[originallemma][0][]0377}{[1][0][]source}{}{}{}}
\makeatother
\title{Sections of a sheaf on a proper scheme over a henselian pair agree with sections on the closed fibre}
\author{The Stacks Project Authors}
\date{stacks-project:etale-cohomology.tex:local-snapshot-20261011; GNU FDL 1.2 or later}
\begin{document}
\maketitle
\noindent\textit{Editorial note.} All delivered mathematical text is the author's own, copied from the pinned Stacks Project chapter file etale-cohomology.tex; no statement or proof step was rewritten, reordered or removed.\par\medskip
\noindent\textit{Editorial note.} The counted proof is the author's complete argument: the proof is written out in full: the topological criterion, the reduction to irreducible X and to the Stein factorisation, the connectedness of the intersection of the closed image with V(I), and the final connected-fibres argument.\par\medskip
\noindent\textit{Editorial note.} The label lemma-zariski-h0-proper-over-henselian-pair names the counted statement, and the proof environment immediately following it in the source is delivered verbatim as its proof.\par\medskip
\noindent\textit{Editorial note.} The counted statement is carried by the author's own lemma environment in the chapter source.\par\medskip
\noindent\textit{Editorial note.} The statement and proof are copied from the pinned chapter file; the six cited inputs are kept as external pointers with their canonical Stacks tags.\par\medskip
\noindent\textit{Editorial note.} The preamble is the chapter's own native preamble excerpt already used for this source: the Xy-pic setup (native preamble lines 14-18) and the native theorem declarations with the mathematical macros (native preamble lines 154-217).\par\medskip
\noindent\textit{Allowed standard prerequisites (editorial).}\begin{itemize}
\item The topological h\textasciicircum{}0 criterion via connectedness of the intersections. Same chapter, the lemma reducing the comparison of sections over X and Z to the connectedness of Y cap Z for every irreducible closed subscheme Y of X. Supplies the criterion that the proof verifies.
\item Spectrality of the underlying space of a proper scheme over a henselian pair. Properties, the lemma making the underlying topological space of X spectral. Supplies the first hypothesis of the topological criterion.
\item Stein factorisation and integrality over a henselian pair. More on Morphisms, the Stein factorisation theorem, and More on Algebra, the lemmas that an integral extension of a henselian pair is again henselian and that the image of an irreducible scheme meets V(I) in a connected set. Supplies the reduction to geometrically connected fibres and the connectedness of T cap V(I).
\end{itemize}\medskip
\ifdefined\hypertarget\hypertarget{sourcepointer-1}{}\fi
\noindent\textit{Source reference.} Original reference 09ZG: sources/stacks-etale-cohomology/pinned-original-source.tex, source lines 15099--15099. Prerequisite pointer invoked at proof line(s) 18176; it names a standard fact of the same chapter that the counted argument uses, and every step of the counted proof built on it is retained verbatim in the delivered body.\par\medskip
\ifdefined\hypertarget\hypertarget{sourcepointer-2}{}\fi
\noindent\textit{Source reference.} Original lemma 09XK: sources/stacks-chow/tags.txt, source lines 12533--12533. Prerequisite pointer invoked at proof line(s) 18187; it names a standard fact of a sibling chapter that the counted argument uses, and every step of the counted proof built on it is retained verbatim in the delivered body.\par\medskip
\ifdefined\hypertarget\hypertarget{sourcepointer-3}{}\fi
\noindent\textit{Source reference.} Original lemma 09Y6: sources/stacks-chow/tags.txt, source lines 12554--12554. Prerequisite pointer invoked at proof line(s) 18192; it names a standard fact of a sibling chapter that the counted argument uses, and every step of the counted proof built on it is retained verbatim in the delivered body.\par\medskip
\ifdefined\hypertarget\hypertarget{sourcepointer-4}{}\fi
\noindent\textit{Source reference.} Original theorem 03H2: sources/stacks-chow/tags.txt, source lines 4549--4549. Prerequisite pointer invoked at proof line(s) 18185; it names a standard fact of a sibling chapter that the counted argument uses, and every step of the counted proof built on it is retained verbatim in the delivered body.\par\medskip
\ifdefined\hypertarget\hypertarget{sourcepointer-5}{}\fi
\noindent\textit{Source reference.} Original lemma 094L: sources/stacks-chow/tags.txt, source lines 11554--11554. Prerequisite pointer invoked at proof line(s) 18178; it names a standard fact of a sibling chapter that the counted argument uses, and every step of the counted proof built on it is retained verbatim in the delivered body.\par\medskip
\ifdefined\hypertarget\hypertarget{sourcepointer-6}{}\fi
\noindent\textit{Source reference.} Original lemma 0377: sources/stacks-chow/tags.txt, source lines 4190--4190. Prerequisite pointer invoked at proof line(s) 18196; it names a standard fact of a sibling chapter that the counted argument uses, and every step of the counted proof built on it is retained verbatim in the delivered body.\par\medskip
\begin{lemma}
\label{lemma-zariski-h0-proper-over-henselian-pair}
Let $(A, I)$ be a henselian pair. Let $f : X \to \Spec(A)$ be a proper morphism
of schemes. Let $Z = X \times_{\Spec(A)} \Spec(A/I)$. For any
sheaf $\mathcal{F}$ on the topological space associated to $X$ we
have $\Gamma(X, \mathcal{F}) = \Gamma(Z, \mathcal{F}|_Z)$.
\end{lemma}

\begin{proof}
We will use Lemma \ref{lemma-h0-topological} to prove this. First observe
that the underlying topological space of $X$ is spectral by Properties, Lemma
\ref{properties-lemma-quasi-compact-quasi-separated-spectral}.
Let $Y \subset X$ be an irreducible closed subscheme. To finish the proof
we show that $Y \cap Z = Y \times_{\Spec(A)} \Spec(A/I)$ is connected.
Replacing $X$ by $Y$
we may assume that $X$ is irreducible and we have to show that $Z$
is connected. Let $X \to \Spec(B) \to \Spec(A)$ be the Stein factorization
of $f$ (More on Morphisms, Theorem
\ref{more-morphisms-theorem-stein-factorization-general}).
Then $A \to B$ is integral and $(B, IB)$ is a henselian pair
(More on Algebra, Lemma \ref{more-algebra-lemma-integral-over-henselian-pair}).
Thus we may assume the fibres of $X \to \Spec(A)$ are geometrically
connected. On the other hand, the image $T \subset \Spec(A)$ of $f$
is irreducible and closed as $X$ is proper over $A$. Hence $T \cap V(I)$
is connected by More on Algebra, Lemma
\ref{more-algebra-lemma-irreducible-henselian-pair-connected}.
Now $Y \times_{\Spec(A)} \Spec(A/I) \to T \cap V(I)$
is a surjective closed map with connected fibres.
The result now follows from Topology, Lemma
\ref{topology-lemma-connected-fibres-quotient-topology-connected-components}.
\end{proof}

\end{document}
